\documentclass[11pt]{article}
\usepackage[margin=1in]{geometry}
\usepackage[T1]{fontenc}
\usepackage{lmodern}
\usepackage{microtype}
\usepackage{setspace}
\usepackage[colorlinks=true,urlcolor=black]{hyperref}
\setlength{\parindent}{0pt}
\setlength{\parskip}{8pt}
\linespread{1.08}
\pagestyle{plain}

\begin{document}

\begin{center}
{\large\bfseries Dimitri Chrysafis\par}
\vspace{2mm}
{\normalsize Statement of Purpose, ETH Zurich MSc Computer Science\par}
\end{center}

\vspace{4mm}

I am a computer science student at the University of Wisconsin--Madison, where I expect to finish my bachelor's degree in May 2027, two years after starting. The compressed timeline has forced me to choose carefully what I work on. Most of that work has centered on the same problem: how AI systems behave when they depend on imperfect data, limited memory, and evaluation methods that can themselves be wrong. I want to study that problem in a more rigorous setting.

My interest in data systems became concrete through a project I built around my own cycling records. While training for Ironman Wisconsin, I accumulated months of Zwift rides spread across interrupted sessions, corrupted FIT files, and short fragments that looked like real workouts. I wanted to answer simple questions about training, but first I had to decide what counted as a ride and which records could be trusted. I built a pipeline that repairs incomplete binary headers, preserves decode errors, stores raw documents in MongoDB, curates sessions in MySQL, and indexes text and embeddings in Elasticsearch. The current archive groups 49 activity files into 31 sessions with 143,400 timestamped records. The project started as a personal tool, but it taught me a broader lesson: an AI application can produce a confident answer while silently using the wrong evidence.

A second project pushed me toward inference systems. I built a local engine for a 125B-parameter sparse mixture-of-experts Qwen model on a 36 GB laptop by streaming quantized experts from NVMe storage. Because the machine could not hold the model in memory, decoding speed depended on which experts were loaded before they were needed. I trained small per-layer predictors from routing traces, added prefetching, and found that the cache also had to pin prefetched experts until use. Across eight held-out prompts, median expert-read wait fell from 37.9 seconds to 13.3 seconds, while the generated tokens matched the baseline. The mixed throughput result, with improvement on five prompts and regression on three, convinced me that inference behavior cannot be understood from a single benchmark number.

I care about evaluation for the same reason. In spackle, a Rust coding agent I wrote for local models, I found that an agent could appear to fix a bug by editing or deleting the tests used to measure success. Adding protected-test comparison changed the measured solve rate of one cloud model from 40\% to 19\%. My open-source work follows the same pattern. My merged pull requests in Unsloth and mlx-audio mostly repair small but important correctness issues, such as preserving system messages, decoding datasets properly, applying sampler parameters, and reporting failures honestly. I enjoy this work because the details determine whether a system is merely impressive or actually dependable.

I want to earn a master's degree now because I have reached the limit of what I can test alone. Finishing my undergraduate degree quickly gave me room for independent projects, but not the sustained research mentorship I need to evaluate whether these design choices generalize. I want advanced coursework in data management, machine learning, optimization, and systems, followed by a substantial thesis. My longer-term goal is to build infrastructure for large models and AI agents across retrieval, inference, and evaluation.

ETH Zurich is a strong fit because its computer science master's program lets me combine the Data Management Systems major with machine intelligence coursework and a full research thesis. Courses such as Data Management Systems, Big Data, Cloud Computing Architecture, Advanced Machine Learning, Optimization for Data Science, and Large Language Models map directly onto the questions I am already investigating. The program's emphasis on independent research is especially important to me. I want a setting where my workloads, baselines, and assumptions are challenged, and where a thesis can turn the questions I began in local projects into defensible research. ETH offers the combination of systems depth, mathematical rigor, and research intensity that I am looking for.

\end{document}
